\documentclass[letterpaper,11pt]{article}
\usepackage[top=0.8in, bottom=0.8in, left=0.9in, right=0.9in]{geometry}
\usepackage[hidelinks]{hyperref}
\usepackage{lmodern}
\usepackage{parskip}
\usepackage{microtype}
\setlength{\parskip}{6pt}
\setlength{\parindent}{0pt}
\begin{document}
\begin{flushleft}
\textbf{\large Ajay Venkatesh} \\
Brooklyn, NY \\
+1 516 585 9013 \\
\href{mailto:av3855@nyu.edu}{av3855@nyu.edu}
\end{flushleft}
\vspace{10pt}
May 21, 2026
\vspace{6pt}

Hiring Manager \\
T-Mobile

\vspace{10pt}

Dear Hiring Manager,

My name is Ajay Venkatesh. I'm finishing my M.S. in Computer Engineering at NYU, with production software experience at PwC in Bengaluru, a summer SDE internship at Amazon, and ongoing on-campus work at NYU's Novel AI Technologies. I'm writing about the Data Scientist role on the Modeling and Valuation team in Finance.

The role's framing of operating and improving financial liability models maps cleanly onto where I've spent a lot of my engineering time. My Big Data graduate project built a \textbf{PySpark / HDFS} pipeline over millions of flight records, training Random Forest and Gradient Boosted Tree models with rigorous EDA and feature engineering. I've also worked on \textbf{RoBERTa fine-tuning with custom LoRA} and \textbf{ResNet} architectures under strict parameter budgets, all of which builds the discipline of producing models that hold up under stakeholder scrutiny.

On production data analytics: at PwC I spent the past few years on a distributed backend integrating with \textbf{Microsoft CRM} over \textbf{SQL Server}, writing complex queries, stored procedures, and parallelized batch ingestion that materially cut runtimes. I built \textbf{Power BI} dashboards over multi-source datasets. At NYU I engineered \textbf{ETL pipelines} transforming \textbf{NoSQL} into \textbf{PostgreSQL} with concurrency control, then powered \textbf{AWS QuickSight} dashboards for real-time analytics, exactly the kind of insight-delivery work the JD describes.

On stakeholder coordination: I've owned features end-to-end alongside architects, product managers, and business stakeholders at PwC, translating ambiguous requirements into reporting and analytical deliverables for non-technical audiences.

Stack-wise: \textbf{SQL}, \textbf{Python}, \textbf{Pandas}, \textbf{NumPy}, \textbf{PySpark}, plus \textbf{scikit-learn} and \textbf{PyTorch} for modeling. Comfortable with \textbf{Power BI}, \textbf{Tableau}, \textbf{QuickSight} for reporting, and \textbf{Snowflake} for data warehousing.

What pulls me toward T-Mobile Finance is the discipline of the work. Liability modeling and valuation forecasts run on accuracy that real financial decisions depend on, where careful analysis matters more than dashboard novelty.

I'd welcome the chance to discuss further. Thanks for considering my application.

\vspace{12pt}
Sincerely, \\
Ajay Venkatesh

\end{document}